\documentclass[11pt, answers]{exam}
\usepackage{tikz}
\usetikzlibrary{decorations.pathreplacing}
\usepackage{amsmath}
\usepackage{amsthm}
\usepackage[boxed]{algorithm}
\usepackage[noend]{algpseudocode}
\usepackage{babel}
\usepackage{titlesec}
\usepackage{changepage}
\usepackage{amsfonts}
\usepackage{amssymb}
\usepackage{commath}
\usepackage{mathrsfs}
\usepackage{fancybox}
\usepackage{xcolor}
\usepackage{enumitem}
\usepackage[noend]{algpseudocode}
\usepackage[margin = 1in]{geometry}
\usepackage[utf8]{inputenc}
\usepackage{pgfplots}
\pgfplotsset{compat=1.18}
\usepackage{graphicx}
\usetikzlibrary{calc}
\newif\ifsteps
\stepstrue
%\stepsfalse

\renewcommand{\solutiontitle}{\noindent}
\renewcommand{\qedsymbol}{$\blacksquare$}
\newcommand{\x}{\times}
\newcommand{\ra}{\rightarrow}
\newcommand{\imp}{\implies}
\newcommand{\ex}{\exists}
\newcommand{\prob}{\textbf{P}}
\newcommand{\R}{\mathbb{R}}
\newcommand{\C}{\mathbb{C}}
\newcommand{\F}{\mathbb{F}}
\newcommand{\N}{\mathbb{N}}
\newcommand{\Z}{\mathbb{Z}}
\newcommand{\Q}{\mathbb{Q}}
\newcommand{\comp}{\mathsf{c}}
\newcommand{\topo}{\mathcal{T}}
\newcommand{\vter}{\mathit{v_{\text{ter}}}}
\newcommand{\vex}{\mathit{v_{\text{ex}}}}
\newcommand{\lang}{\mathcal{L}}
\newcommand{\ham}{\mathcal{H}}
\newcommand{\ess}{\mathcal{S}}
\newcommand{\avg}[1]{\langle {#1} \rangle}
\newcommand{\bigavg}[1]{\left\langle {#1} \right\rangle}
\newcommand{\allint}{\int_{-\infty}^{+\infty}}
\newcommand\todo[1]{\textcolor{red}{#1}}
\newcommand{\fullwidthseed}{\par\noindent\rule{\linewidth}{0pt}\vspace*{-\topskip}}
\newcommand{\vect}[1]{|{#1}\rangle}
\newcommand{\schro}{Schr{\"o}dinger}
\allowdisplaybreaks
\usepackage{calligra}
\DeclareMathAlphabet{\mathcalligra}{T1}{calligra}{m}{n}
\DeclareFontShape{T1}{calligra}{m}{n}{<->s*[2.2]callig15}{}
\newcommand{\scriptr}{\mathcalligra{r}\,}
\newcommand{\boldscriptr}{\pmb{\mathcalligra{r}}\,}
\usepackage{comment}
\includecomment{steps}
\usepackage{braket}

\begin{document}
	\begin{center}
		\Large \bf Physics 7501 Quiz 1 \\ 9/4/2026 \\ Sharon Ye
	\end{center}
	
	\section*{(1.1.3) Pauli matrices}
	
	Recall the Pauli matrices
	\begin{align*}
		\hat{X}
		&= \begin{pmatrix}
			0 & 1 \\
			1 & 0
		\end{pmatrix},
		\qquad
		\hat{Y}
		= \begin{pmatrix}
			0 & -i \\
			i & 0
		\end{pmatrix},
		\qquad
		\hat{Z}
		= \begin{pmatrix}
			1 & 0 \\
			0 & -1
		\end{pmatrix}.
		\tag{1.1}
	\end{align*}
	Express the following matrix $\hat{M}$ as a linear superposition of
	identity matrix $\hat{1}_{2\times 2}$ and Pauli matrices:
	\begin{align*}
		\hat{M}
		&= \begin{pmatrix}
			3 & 6 \\
			4 & 5
		\end{pmatrix}
		= m_0\hat{1}_{2\times 2}
		+ m_1\hat{X}
		+ m_2\hat{Y}
		+ m_3\hat{Z}.
		\tag{1.2}
	\end{align*}
	Obtain $m_{0,1,2,3}$.
	\begin{solution}
		For $\hat{\sigma}_0=\hat{1},\hat{\sigma}_1=\hat{X},\hat{\sigma}_2=\hat{Y},\hat{\sigma}_3=\hat{Z}$, note that under the Hilbert-Schmidt norm $\braket{\psi_i|\psi_j}=\operatorname{tr}(\psi_i^\dagger\psi_j)=2\delta_{ij}$. Thus $m_i = \frac{1}{2}\braket{\hat{\sigma}_i|\hat{M}} \implies$
		\begin{align*}
			m_0 &= \frac{1}{2}\operatorname{tr}
			\left[\begin{pmatrix}
				1 & 0 \\
				0 & 1
			\end{pmatrix}
			\begin{pmatrix}
				3 & 6 \\
				4 & 5
			\end{pmatrix}\right] = 4 \\
			m_1 & = \frac{1}{2}\operatorname{tr}
			\left[\begin{pmatrix}
				0 & 1 \\
				1 & 0
			\end{pmatrix}
			\begin{pmatrix}
				3 & 6 \\
				4 & 5
			\end{pmatrix}\right] = 5 \\
			m_2 &= \frac{1}{2}\operatorname{tr}
			\left[\begin{pmatrix}
				0 & -i \\
				i & 0
			\end{pmatrix}
			\begin{pmatrix}
				3 & 6 \\
				4 & 5
			\end{pmatrix}\right]
			= i \\
			m_3 &= \frac{1}{2}
			\operatorname{tr}
			\left[\begin{pmatrix}
				1 & 0 \\
				0 & -1
			\end{pmatrix}
			\begin{pmatrix}
				3 & 6 \\
				4 & 5
			\end{pmatrix}\right] = -1
		\end{align*}
	\end{solution}
\end{document}
